% Generated draft: requires mathematical and source review.
\documentclass[11pt]{article}
\usepackage[margin=25mm]{geometry}
\usepackage{fontspec}
\setmainfont{Latin Modern Roman}
\usepackage{amsmath,amssymb,mathtools,unicode-math}
\setmathfont{Latin Modern Math}
\usepackage{enumitem,needspace,xurl,hyperref}
\setlength{\parindent}{0pt}
\setlength{\parskip}{5pt}
\setlength{\emergencystretch}{3em}
\newcommand{\N}{\mathbb{N}}
\newcommand{\Z}{\mathbb{Z}}
\newcommand{\Q}{\mathbb{Q}}
\newcommand{\R}{\mathbb{R}}
\newcommand{\C}{\mathbb{C}}
\newcommand{\F}{\mathbb{F}}
\newcommand{\A}{\mathbb{A}}
\newcommand{\PP}{\mathbb{P}}
\newcommand{\E}{\mathbb{E}}
\newcommand{\eps}{\varepsilon}
\newcommand{\dd}{\,\mathrm{d}}
\begin{document}
\section*{UnsolvedMath ID: 30002863}
\textbf{Source problem code:} catalogue entry
\subsection{Definitions and mathematical statement}
The supplied record poses three classification questions, but does not specify all the data needed for a unique mathematical formulation.

For the elementary meaning of linear precision, let \(d\geq 1\), let \(A\subset\R^d\) be a finite set, and let \(D\subset\R^d\) be a parameter domain. Suppose that functions \(\lambda_a:D\to\R\), indexed by \(a\in A\), satisfy
\[
\sum_{a\in A}\lambda_a(x)=1\qquad(x\in D).
\]
Such functions have \emph{linear precision} if
\[
\sum_{a\in A}a\,\lambda_a(x)=x\qquad(x\in D).
\]
Equivalently, they reproduce every affine function:
\[
\forall c\in\R\;\forall v\in\R^d\;\forall x\in D,
\qquad
\sum_{a\in A}(c+v\cdot a)\lambda_a(x)=c+v\cdot x.
\]
This explains the affine reproduction property; it does not supply the missing definition of the toric blending functions or the admissible domains in the source question. In particular, the source does not specify whether a precision identity must hold on the entire patch domain, including its boundary, or only where the defining expressions are defined.

For a field \(k\), a birational self-map of projective three-space is a rational map
\[
f:\PP^3_k\dashrightarrow\PP^3_k
\]
for which there is a rational map \(g:\PP^3_k\dashrightarrow\PP^3_k\) satisfying
\[
g\circ f=\operatorname{id}_{\PP^3_k},\qquad
f\circ g=\operatorname{id}_{\PP^3_k}
\]
as rational maps, hence on appropriate dense open subsets. Here \(k\) is a placeholder: the source does not choose a field or its characteristic.

\textbf{Requested classifications.} The supplied statement asks:
\begin{enumerate}
\item Which toric Bezier patches have linear precision?
\item Which toric Bezier volumes have rational linear precision?
\item How can birational maps of projective \(3\)-space be classified?
\end{enumerate}
The source does not define the admissible toric patches and volumes, their lattice configurations, weights, or parameterizations. It also does not define rational linear precision or specify which rational changes of parameter are allowed. Consequently, rational linear precision is retained as an unresolved source term; it must not be replaced merely by rationality of the blending functions. No equivalence relation or other classification criterion is specified for the third question, and no precise mathematical relationship among the three requested classifications is supplied. These omissions prevent a more specific unconditional classification statement.
\subsection{Short English statement}
The problem asks for classifications of toric Bezier patches with linear precision, toric Bezier volumes with rational linear precision, and birational maps of projective three-space. Linear precision means reproducing affine coordinates through the blending functions. The admissible toric data, the meaning of rational linear precision, and the field and classification criterion for projective maps require clarification.
\subsection{Sources}
\begin{itemize}
\item Supplied local UnsolvedMath registry record, permanent ID \(30002863\), source code \texttt{catalogue entry}, in the Oberwolfach Reports Open Problems collection. The local registry is \texttt{data/unsolvedmath/problems.json}, with additional existing records maintained in \texttt{catalog/problems.json}. This draft uses the supplied record; it does not claim independent inspection of those files.
\item Supplied source citation: \emph{Multivariate Splines and Algebraic Geometry (2016)}, Oberwolfach Reports, DOI 10.4171/owr/2015/21, pp.\ 1139--1200. \url{https://doi.org/10.4171/owr/2015/21}. The literature-source list labels this reference ``Oberwolfach Report 21.'' The record gives proposed year 2015 and names no individual proposer. The relationship between that year and the citation's parenthetical 2016 requires verification.
\item Supplied literature reference: \emph{Linear Precision for Toric Surface Patches}. \url{https://doi.org/10.1007/s10208-009-9052-6}.
\item Supplied literature reference: \emph{Families of Polytopes with Rational Linear Precision}. \url{https://doi.org/10.1007/s10208-022-09583-7}.
\item Literature assessment dated \texttt{2026-08-21}: the supplied curation labels the problem partially solved and reports classifications of linear-precision toric surface patches in important cases, known higher-dimensional families with rational linear precision, and an incomplete requested general classification in dimension three. This is supplied dated curation derived from the OWR record, not an independently verified mathematical or current literature assessment. The record also labels its statement exact and retained byte-for-byte; that is a supplied transcription assessment and does not resolve the formulation issues above.
\end{itemize}
\end{document}
